\documentclass[12pt]{article}
% Packages
\usepackage{amsmath,amssymb}
\usepackage{enumitem}
\usepackage{hyperref}

% Author & date
\title{Octonions, \(G_2\), and Number-Theoretic Obstructions in the Arithmetic of Order Framework: \\ 
A Topological–Algebraic Explanation for Electroweak Mass under Reinterpretation of the Higgs Mechanism}
\author{Faysal El Khettabi \\ 
Ensemble AIs \\ 
\texttt{faysal.elkhettabi@gmail.com}}
\date{August 22, 2025}

\begin{document}

\maketitle

\section*{Dedication} 
\begin{center}
\textit{This work is dedicated to our Parents. \\ 
The Arithmetic of Order (AoO) is founded on the natural counting sequence 1 to \(n\) to \(n+1\), a principle we learn from our Mothers (XX) and our Fathers (XY).}
\end{center}

\begin{abstract}
We resolve a sixty-year-old geometric puzzle—the non-existence of a complex structure on the 6-sphere—by revealing its origin in number theory, while simultaneously uncovering a novel mechanism for the origin of electroweak mass. The known geometric obstruction (non-vanishing Nijenhuis tensor) is shown to be a direct consequence of a deeper arithmetic obstruction: the Brauer–Manin obstruction arising from the failure of the equation \(x^2+1=0\) to have solutions over \(\mathbb{Q}_p\) for primes \(p \equiv 3 \pmod{4}\), including the prime \(p=7\). Within the finitistic Arithmetic of Order (AoO) framework, this unified mathematical “no” becomes a physical “yes”: the primes themselves impose intrinsic non-associativity in the octonionic algebra of physical states, which manifests as irreducible geometric torsion on \(S^6(\mathbb{R})\). Consequently, the masses of the \(W\) and \(Z\) bosons are not free parameters but \emph{computable, number-theoretic inevitabilities}, arithmetically mandated by the properties of the primes themselves.
\end{abstract}

\subsection*{Prime \(7\) and Limits of the Standard Complex Structure}

While standard quantum mechanics employs the complex unit \(i\) with \(i^2=-1\) globally, this implicitly assumes a fully integrable complex structure. In the AoO framework, the prime \(p=7\) plays a critical role: the equation \(x^2+1=0\) has no solution in \(\mathbb{Q}_7\), the 7-adic completion of the rationals, reflecting a local arithmetic obstruction. This Brauer–Manin obstruction at \(p=7\) prevents a globally consistent complex structure on the rational 6-sphere \(S^6(\mathbb{Q})\subset \mathbb{Q}^7\). Consequently, massless propagation of electroweak gauge bosons, which would require a global \(i\), is forbidden, and the nonzero \(W\) and \(Z\) masses emerge as a direct, number-theoretically mandated consequence of the primes themselves.

\section{Finitistic Foundations and Emergent Geometry}

The AoO philosophy posits that reality is fundamentally discrete and computable, built upon rational numbers (\(\mathbb{Q}\)) generated constructively rather than assumed as a continuum.

\begin{itemize}
    \item \emph{Rational Octonions via Farey Recursion.} Starting with a finite-yet-unbounded construction of \(\mathbb{Q}\) through iterated Farey sequences, the 8-dimensional octonion algebra over \(\mathbb{Q}\), denoted \(\mathbb{O}_{\mathbb{Q}}\), is defined, along with its 7-dimensional imaginary subspace \(\operatorname{Im}\mathbb{O}_{\mathbb{Q}} \cong \mathbb{Q}^7\).
    \item \emph{The Rational 6-Sphere and its Emergent Continuum.} The unit-norm locus in \(\operatorname{Im}\mathbb{O}_{\mathbb{Q}}\) defines the rational 6-sphere, \(S^6(\mathbb{Q})\). As denominators in the Farey construction grow, this dense set converges metrically to the real 6-sphere \(S^6(\mathbb{R})\), making the continuum an emergent approximation.
    \item \emph{Coset Geometry: The Rosetta Stone.} The full coset \(G_2/(SU(2)_L \times SU(2)_R)\) is an 8-dimensional algebraic variety over \(\mathbb{Q}\). A distinguished 6-dimensional submanifold is isomorphic to \(S^6(\mathbb{Q})\), realized both as the unit sphere in the imaginary octonions and as the coset \(G_2/SU(3)\). This submanifold translates octonionic algebra into physical geometry; the continuum \(S^6(\mathbb{R})\) is its emergent, asymptotic limit.
\end{itemize}

\section{Physical Algebra: States, Associators, and Interaction as Geometry}

In AoO, fundamental states and their interactions are described by non-associative octonionic algebra, directly inducing the emergent geometry.

\begin{itemize}
    \item \emph{Physical States.} Fundamental states are tri-octonion vectors:
    \[
    v = (o_1, o_2, o_3) \in (\operatorname{Im}\mathbb{O}_{\mathbb{Q}})^3.
    \]
    \item \emph{Associator as Internal Tension.} Each state's associator
    \[
    [v] := (o_1 o_2) o_3 - o_1 (o_2 o_3)
    \]
    quantifies intrinsic non-associativity, interpreted as “internal tension” or charge.
    \item \emph{Interaction and Emergent Complex Structure.} Interactions between states \(v_1, v_2\) are governed by the commutator of associators:
    \[
    [[v_1],[v_2]] := [v_1][v_2] - [v_2][v_1] = 2\, [v_2] \times [v_1],
    \]
    where the factor of 2 arises from structure constants. The octonionic cross product defines the almost-complex structure \(J\) on \(S^6(\mathbb{R})\):
    \[
    J_g(v) = g \times v, \quad g \in S^6.
    \]
    \textbf{Thus, physical interaction is the geometric structure itself.}
\end{itemize}

\section{Geometric Obstruction: Non-Integrability from Non-Associativity}

Non-associativity of octonions induces non-integrability of \(J\).

\begin{itemize}
    \item \emph{Nijenhuis Tensor.} An almost-complex structure \(J\) is integrable if and only if
    \[
    N_J(X,Y) = [JX, JY] - J[JX,Y] - J[X,JY] - [X,Y] = 0
    \]
    for all vector fields \(X,Y\).
    \item \emph{Obstruction Manifest.} For \(J_g(v) = g \times v\), \(N_J \neq 0\), classically proving that \(S^6\) has no global integrable complex structure.
    \item \emph{AoO Interpretation.} Non-vanishing \(N_J\) is the macroscopic manifestation of irreducible internal tensions \([v] \neq 0\), which directly correspond to geometric torsion.
\end{itemize}

\section{Arithmetic Obstruction: Brauer–Manin and the Prime \(7\)}

The “why” of non-integrability lies in number theory.

\begin{itemize}
    \item \emph{Local-to-Global Principle.} The Hasse principle states that an equation has a solution in \(\mathbb{Q}\) only if it has solutions over \(\mathbb{R}\) and all \(\mathbb{Q}_p\).
    \item \emph{Quadratic Prototype.} \(x^2 + 1 = 0\) is solvable in \(\mathbb{Q}_p\) iff \(p \equiv 1 \pmod{4}\) or \(p=2\). For \(p \equiv 3 \pmod{4}\) (including \(p=7\)), no solution exists; \(-1\) is not a quadratic residue.
    \item \emph{Brauer–Manin Obstruction.} Extending to the rational octonion algebra \(B\) (where \(B \otimes_{\mathbb{Q}} \mathbb{R} \cong \mathbb{O}\)), the ability to define a complex structure on \(S^6(\mathbb{Q}) \subset \mathbb{Q}^7\) depends on whether \(B \otimes_{\mathbb{Q}} \mathbb{Q}_p\) splits or remains a division algebra:
    \begin{itemize}
        \item Over \(\mathbb{R}\): \(B\) splits; \(J\) exists.
        \item \(p \equiv 1 \pmod{4}\): \(B\) is expected to split; a local complex structure exists.
        \item \textbf{\(p \equiv 3 \pmod{4}\), including \(p=7\)}: \(B\) remains division; \textbf{no local complex structure exists.}
    \end{itemize}
    The absence of local solutions at \(p=7\) and other \(3 \pmod{4}\) primes constitutes a Brauer–Manin obstruction, proving that no global integrable complex structure exists on \(S^6(\mathbb{Q})\). The prime \(7\) is thus intrinsic to the obstruction, providing the 7-adic topology \(\mathbb{Q}_7\) relevant to the 7-dimensional space \(\mathbb{Q}^7\) in which \(S^6(\mathbb{Q})\) is defined.
\end{itemize}

\section{Unified Physical Law: Mass as Arithmetic Inevitability}

The geometric obstruction (\(N_J \neq 0\)) is a consequence of the arithmetic obstruction (Brauer–Manin).

\begin{itemize}
    \item The number-theoretic obstruction \textbf{mandates} non-associativity at the fundamental level \([v] \neq 0\).
    \item Non-associativity \textbf{manifests} as irreducible torsion in \(J\) on the emergent continuum.
    \item This torsion acts as unavoidable resistance to propagation; massless electroweak gauge bosons would require a globally integrable complex structure—explicitly forbidden by the primes, including \(p=7\).
\end{itemize}

\paragraph{Conclusion}
The \(W\) and \(Z\) bosons acquire mass not as a contingent Higgs parameter but as a \emph{computable, number-theoretic inevitability}. This reinterprets the Higgs mechanism as a fundamental arithmetic-geometric constraint on the structure of the universe.

\bigskip
\hrulefill
\bigskip

\noindent\emph{Selected References}
\begin{enumerate}[label={[\arabic*]}]
    \item F. El Khettabi, \emph{Octonions, \(G_2\), and Knot Topology in the Arithmetic of Order Framework}, 2025.\\
    \texttt{https://efaysal.github.io/HCNFEK2024FE/HIGGSMECHANISMSAOOFEKAUGUST2025.pdf}
    \item S. Kobayashi, K. Nomizu, \emph{Foundations of Differential Geometry, Vol.~II}, Wiley Classics Library, 1969.
\end{enumerate}

\end{document}

